\documentclass[10pt,a4paper]{amsart}
\usepackage[margin=19mm]{geometry}
\usepackage{amsmath,amssymb,mathtools}
\usepackage[scr=rsfs,cal=euler]{mathalfa}
\usepackage{xfrac}
\usepackage[unicode,hidelinks,bookmarks=false]{hyperref}
\newcommand{\ie}{i.e.\ }
\newcommand{\cf}{cf.\ }
\newcommand{\resp}{respectively\ }
\newcommand{\Subsection}{\subsection}
\providecommand{\nobreakdash}{\nobreak\hbox{-}}
\setlength{\emergencystretch}{2em}
\multlinegap0pt
\usepackage{tikz}
\usetikzlibrary{cd}
\usepackage{mathtools}
\usepackage{bm}
\newtheorem{lemma}{Lemma}
\usepackage{cleveref}
\crefname{lemma}{Lemma}{Lemmas}
\newcommand{\D}{{\mathbb D}}
\newcommand\cW{{\mathcal W}}
\newcommand{\cY}{\mathcal{Y}}
\title{Asymptotics of small eigenvalues on degenerations of K\"ahler manifolds}
\author{Junyu Cao}
\date{Source: arXiv:2605.08023v1 (CC BY 4.0)}
\begin{document}
\maketitle

\noindent\emph{Source and licence.} Asymptotics of small eigenvalues on degenerations of K\"ahler manifolds. Version arXiv:2605.08023v1 (CC BY 4.0), distributed by arXiv under the Creative Commons Attribution 4.0 International licence, http://creativecommons.org/licenses/by/4.0/, as recorded in the arXiv licence metadata for this exact version and shown on the abstract page https://arxiv.org/abs/2605.08023v1. Full licence notice: \texttt{licenses/arxiv-2605.08023-CC-BY-4.0.txt}.\par\smallskip

\noindent\textit{Editorial scope.} Counted unit: the article's lemma lemma (source0.tex lines 985--1017). The article's complete original proof of it is delivered with it (source0.tex lines 1019--1115, one \texttt{proof} environment of 97 lines). No mathematics is written, repaired or re-derived: the statement and the proof are the article's own lines, in the article's own notation, and no retained line is substituted or re-flowed.  Retained uncounted as the source-local context the proof needs: lines 965--974.  Nothing was omitted from any retained span, so the package carries no omission record.  No retained line contains a citation, so the wrapper carries no bibliography and no bibliography entry.  The retained lines contain the article's own diagram code, which the wrapper typesets with the standard tikz packages; no image file is delivered and the package declares no asset dependency.  Editorial formatting only, in addition: an AMS \texttt{amsart} preamble that restates the article's own declarations and macros used by the retained lines, namely the environments \texttt{lemma} on separate counters, together with the standard packages those lines need.  The retained environments print in this document as Lemma 1 in order of appearance, whereas the article prints the same items with the numbering of its own full document.  The complete native source of the article as distributed in the arXiv e-print for this exact version is bundled unchanged as \texttt{sources/200127/source0.tex}.
\begin{equation}
    \label{diagram:diagram-Z}
    \begin{tikzcd}
    \widehat{F^{-1}X} \arrow[r, "\iota"] \arrow[dr, bend right=20, "\widehat{\Pi}"']
    & F^{-1}X \coloneqq X \times_{\D_s} \D_t \arrow[r, "F"] \arrow[d, "\Pi"]
    & X \arrow[d, "\pi"] \\
    & \D_t \arrow[r, "t \mapsto t^m"']
    & \D_s
\end{tikzcd}
\end{equation}

\begin{lemma}
    \label[lemma]{lem:counting-irrd-cpnt-of-Z}
    For any semistable reduction \(p \colon \cY \to \D_t\) of \(\pi\)
    over a finite base change, fitting into the commutative diagram
    \[
    \begin{tikzcd}
        & \cY  \arrow[r, "\mu"] \arrow[d, "p"] & X \arrow[d, "\pi"] \\
        & \D_t \arrow[r, "f_d"] & \D_s
    \end{tikzcd}
    \]
    where \(f_d(t)=t^d=s\), and the induced map
    \((\mu,p)\colon \cY\to X\times_{\D_s}\D_t\) is an isomorphism over
    \(\D_t^\circ\), we define the following invariant of the
    degeneration:
    \begin{align*}
        N(X)
        &\coloneqq N(\mu^*\omega_X \eqqcolon \beta_\cY,\cY) \\
        &= \#\left\{
        D \subseteq \cY_0
        \;\middle|\;
        D \text{ is an irreducible component and }
        \int_D \beta_\cY^n > 0
        \right\}.
    \end{align*}

    Then
    \begin{itemize}
        \item \(N(X)\) is independent of the choice of semistable
        reduction,
        \item \(N(X)\) is equal to the number of irreducible components
        of \(Z\), which is denoted by \(N\).
    \end{itemize}
\end{lemma}

\begin{proof}
We first show the birational invariance. Let
\(\mu_i\colon \cY_i\to X\), \(i=1,2\), be two
semistable reductions over the same base change \(s=t^d\), and set
\(\beta_i=\mu_i^*\omega_X\). Let \(p_i\colon \cY_i\to\D_t\) be the structure maps.

Since the induced maps \((\mu_i,p_i)\colon \cY_i\to X\times_{\D_s}\D_t\)
are isomorphisms over \(\D_t^\circ\), the varieties \(\cY_1\) and
\(\cY_2\) are birational over \(X\).
We can
choose a common log resolution of pairs $(\cY_1, (\cY_1)_0)$ and $(\cY_2,(\cY_2)_0)$ over \(X\), namely, the following diagram:
\[
\begin{tikzcd}
& \cW \arrow[dl, "\phi_1"'] \arrow[dr, "\phi_2"] & \\
\cY_1 && \cY_2 .
\end{tikzcd}
\]
Then \(\phi_1^*\beta_1=\phi_2^*\beta_2\). For an irreducible component
\(G\subset \cW_0\), either \(G\) is exceptional over \(\cY_i\), in which
case \(\dim \phi_i(G)<n\) and the projection formula gives
\[
\int_G \phi_i^*\beta_i^n=0,
\]
or \(G\) is the strict transform of a unique irreducible component
\(D\subset (\cY_i)_0\), in which case
\[
\int_G \phi_i^*\beta_i^n=\int_D \beta_i^n .
\]
Thus the components with positive \(\beta\)-volume are preserved under
passing to a common log resolution. Hence
\[
N(\beta_1,\cY_1)=N(\beta_2,\cY_2)
\]
for semistable reductions over the same base change.

The same argument shows invariance under further ramified base change.
Indeed, let \(p\colon\cY\to\D_t\) be semistable and pull it back by
\(u\mapsto u^q\). If \(p'\colon\cY'\to\D_u\) is a semistable reduction
of \(\cY\times_{\D_t}\D_u\), and \(\psi\colon\cY'\to\cY\) is the induced map,
then \(\beta_{\cY'}=\psi^*\beta_\cY\). Every component of \((\cY')_0\) is
either the strict transform of a component of \(\cY_0\), or is exceptional
over an intersection stratum of \(\cY_0\). The exceptional components have
zero \(\beta\)-volume by the projection formula, while strict transforms
have the same volume as the original components. Therefore
\[
N(\beta_{\cY'},\cY')=N(\beta_\cY,\cY).
\]
Given two arbitrary semistable reductions of degrees \(d_1\) and \(d_2\),
we pass both to the common base change of degree
\(\ell=\operatorname{lcm}(d_1,d_2)\). The preceding paragraph and the
same-base-change invariance then imply that \(N(X)\) is independent of
the semistable reduction.

It remains to identify this number with the number of irreducible
components of \(Z\). Write
\[
Z=\widehat{\Pi}^{-1}(0)=\bigcup_{j=1}^N Z_j .
\]
After a further base change \(u\mapsto t=u^q\), take a semistable
reduction
\[
\nu\colon \cY\to \widehat{F^{-1}X}\times_{\D_t}\D_u .
\]
Let \(r\colon\cY\to\widehat{F^{-1}X}\) be the induced map and
\(h \coloneqq F \circ \iota \colon\widehat{F^{-1}X}\to X\) be the composition of maps appearing in \cref{diagram:diagram-Z}. Taking \(\mu=h\circ r\), we have
\[
\beta_\cY=\mu^*\omega_X=r^*h^*\omega_X,
\]
and we compute $N(X)$ using this semistable model,
\[
N(X)=N(\beta_\cY,\cY).
\]

Since \(\widehat{F^{-1}X}\) is normal and \(Z\) is reduced, the pullback
family is smooth at the generic point of each \(Z_j\). Hence the
irreducible components of \(\cY_0\) are precisely the strict transforms
\(\widetilde Z_j\), together with \(\nu\)-exceptional divisors. The latter
map to subsets of dimension \(<n\), so they have zero \(\beta_\cY\)-volume.
Since \(h\) is finite over the central fiber, for each \(j\), the map
\[
h_j:=h|_{Z_j}\colon Z_j\to X_0
\]
has image an irreducible component \(D_{\alpha(j)}\subset X_0\) and is
generically finite of degree \(e_j\ge 1\). Therefore
\[
\int_{\widetilde Z_j}\beta_\cY^n
=
e_j\int_{D_{\alpha(j)}}\omega_X^n
>0.
\]
Thus the components of \(\cY_0\) with positive \(\beta_\cY\)-volume are
exactly \(\widetilde Z_1,\dots,\widetilde Z_N\). Consequently
\[
N(X)=N(\beta_\cY,\cY)=N,
\]
as claimed.
\end{proof}

\end{document}
